% =============================================================================
% منحنی‌های آموزش فروریزش خاموش — خروجی tools/tensorboard_to_thesis.py
%
% فایل‌های لازم (هر دو در مخزن نگهداری می‌شوند، پس شکل بدون سرور بازساختنی است):
%   thesis/figures/curves/train_avg_loss_per_epoch.csv
%   thesis/figures/curves/eval_abs_rel.csv
%
% ستون‌های هر CSV: step collapsed collapsed_control fixed fixed_control
%   collapsed  = اجرای هدف عمق مستقیم (فروریخته)
%   fixed      = اجرای هدف عمق معکوس (اصلاح‌شده)
% ستون‌های *_control بازوی فاقد K همان دو اجرا هستند و در شکل رسم نشده‌اند،
% چون در اجرای فروریخته با ستون متناظر خود بیت‌به‌بیت یکسان‌اند.
%
% محل درج: بخش ۴--۴ (پیش از جدول نتیجهٔ ViT-L)، با دستور:
%   % =============================================================================
% منحنی‌های آموزش فروریزش خاموش — خروجی tools/tensorboard_to_thesis.py
%
% فایل‌های لازم (هر دو در مخزن نگهداری می‌شوند، پس شکل بدون سرور بازساختنی است):
%   thesis/figures/curves/train_avg_loss_per_epoch.csv
%   thesis/figures/curves/eval_abs_rel.csv
%
% ستون‌های هر CSV: step collapsed collapsed_control fixed fixed_control
%   collapsed  = اجرای هدف عمق مستقیم (فروریخته)
%   fixed      = اجرای هدف عمق معکوس (اصلاح‌شده)
% ستون‌های *_control بازوی فاقد K همان دو اجرا هستند و در شکل رسم نشده‌اند،
% چون در اجرای فروریخته با ستون متناظر خود بیت‌به‌بیت یکسان‌اند.
%
% محل درج: بخش ۴--۴ (پیش از جدول نتیجهٔ ViT-L)، با دستور:
%   % =============================================================================
% منحنی‌های آموزش فروریزش خاموش — خروجی tools/tensorboard_to_thesis.py
%
% فایل‌های لازم (هر دو در مخزن نگهداری می‌شوند، پس شکل بدون سرور بازساختنی است):
%   thesis/figures/curves/train_avg_loss_per_epoch.csv
%   thesis/figures/curves/eval_abs_rel.csv
%
% ستون‌های هر CSV: step collapsed collapsed_control fixed fixed_control
%   collapsed  = اجرای هدف عمق مستقیم (فروریخته)
%   fixed      = اجرای هدف عمق معکوس (اصلاح‌شده)
% ستون‌های *_control بازوی فاقد K همان دو اجرا هستند و در شکل رسم نشده‌اند،
% چون در اجرای فروریخته با ستون متناظر خود بیت‌به‌بیت یکسان‌اند.
%
% محل درج: بخش ۴--۴ (پیش از جدول نتیجهٔ ViT-L)، با دستور:
%   % =============================================================================
% منحنی‌های آموزش فروریزش خاموش — خروجی tools/tensorboard_to_thesis.py
%
% فایل‌های لازم (هر دو در مخزن نگهداری می‌شوند، پس شکل بدون سرور بازساختنی است):
%   thesis/figures/curves/train_avg_loss_per_epoch.csv
%   thesis/figures/curves/eval_abs_rel.csv
%
% ستون‌های هر CSV: step collapsed collapsed_control fixed fixed_control
%   collapsed  = اجرای هدف عمق مستقیم (فروریخته)
%   fixed      = اجرای هدف عمق معکوس (اصلاح‌شده)
% ستون‌های *_control بازوی فاقد K همان دو اجرا هستند و در شکل رسم نشده‌اند،
% چون در اجرای فروریخته با ستون متناظر خود بیت‌به‌بیت یکسان‌اند.
%
% محل درج: بخش ۴--۴ (پیش از جدول نتیجهٔ ViT-L)، با دستور:
%   \input{./figures/pending/collapse_curves}
% =============================================================================

\begin{figure}[htbp]
    \centering
    \begin{tikzpicture}
    \begin{groupplot}[
        group style={group size=1 by 2, vertical sep=1.05cm, x descriptions at=edge bottom},
        thesisplot,
        width=0.86\textwidth, height=5.0cm,
        xmin=-1, xmax=40,
    ]

    % منحنی‌ها مستقیماً برچسب‌گذاری می‌شوند و راهنمای جعبه‌ای حذف شده است، چون
    % در این دو نمودار جعبهٔ راهنما ناگزیر روی خودِ منحنی‌ها می‌افتد.
    \nextgroupplot[
        ymode=log,
        ymin=0.005, ymax=1.0,
        ylabel={\rl{هزینهٔ آموزش هر دوره}},
    ]
    \addplot[draw=thControl, mark=none] table[x=step, y=collapsed, col sep=space]
        {figures/curves/train_avg_loss_per_epoch.csv};
    \addplot[draw=thIdeal, mark=none] table[x=step, y=fixed, col sep=space]
        {figures/curves/train_avg_loss_per_epoch.csv};
    \node[thnote, anchor=west, text=thControl] at (axis cs:1.5,0.155)
        {\rl{اجرای فروریخته --- هدف: عمق مستقیم}};
    \node[thnote, anchor=west, text=thIdeal] at (axis cs:1.5,0.031)
        {\rl{اجرای اصلاح‌شده --- هدف: عمق معکوس}};

    \nextgroupplot[
        ymin=0, ymax=0.46,
        ylabel={\rl{\lr{AbsRel} اعتبارسنجی}},
        xlabel={\rl{دورهٔ آموزش}},
    ]
    \addplot[draw=thControl, mark=none] table[x=step, y=collapsed, col sep=space]
        {figures/curves/eval_abs_rel.csv};
    \addplot[draw=thIdeal, mark=none] table[x=step, y=fixed, col sep=space]
        {figures/curves/eval_abs_rel.csv};
    \node[thnote, anchor=west, text=thControl] at (axis cs:1.5,0.415)
        {\rl{اجرای فروریخته --- منجمد روی $0.3465$ از دورهٔ نخست}};
    \node[thnote, anchor=west, text=thIdeal] at (axis cs:1.5,0.135)
        {\rl{اجرای اصلاح‌شده --- نزول عادی تا $0.055$}};

    \end{groupplot}
    \end{tikzpicture}
    \caption[امضای فروریزش خاموش در منحنی‌های آموزش]{امضای \emph{فروریزش
    خاموش} توضیح‌داده‌شده در بخش~\ref{subsec:disparity_target_space}، در دو
    اجرای کامل ۴۰ دوره‌ای \lr{ViT-L} که تنها در فضای هدف تفاوت دارند.
    \textbf{بالا:} هزینهٔ آموزش اجرای فروریخته (قرمز) هموار و نزولی است ---
    از \fadecimal{۰٫۴۴۲} به \fadecimal{۰٫۴۲۹} --- و اگر تنها همین منحنی پایش
    شود، آموزش کاملاً سالم به نظر می‌رسد؛ زیرا با خروجی ثابت، مقدار هزینه
    صرفاً تابع آماره‌های هر دسته است. \textbf{پایین:} معیار اعتبارسنجی همان
    اجرا از \emph{دورهٔ نخست} روی \fadecimal{۰٫۳۴۶۵} منجمد می‌ماند و تا پایان
    آموزش حتی در رقم چهارم اعشار تغییر نمی‌کند. نکتهٔ گویا آنکه در این حالت،
    مقادیر اعتبارسنجی دو بازوی \proposedmodel{} و \controlmodel{}
    \emph{بیت‌به‌بیت} یکسان است، چون خروجی هر دو شبکه ثابت صفر شده و دیگر به
    ورودی --- و در نتیجه به $K$ --- وابسته نیست. پس از انتقال هدف به فضای عمق
    معکوس (سبز)، هر دو منحنی رفتار عادی می‌یابند. این شکل نشان می‌دهد که چرا
    پایش هزینهٔ آموزش به‌تنهایی برای تشخیص سلامت آموزش کافی نیست.}
    \label{fig:silent_collapse_curves}
\end{figure}

% =============================================================================

\begin{figure}[htbp]
    \centering
    \begin{tikzpicture}
    \begin{groupplot}[
        group style={group size=1 by 2, vertical sep=1.05cm, x descriptions at=edge bottom},
        thesisplot,
        width=0.86\textwidth, height=5.0cm,
        xmin=-1, xmax=40,
    ]

    % منحنی‌ها مستقیماً برچسب‌گذاری می‌شوند و راهنمای جعبه‌ای حذف شده است، چون
    % در این دو نمودار جعبهٔ راهنما ناگزیر روی خودِ منحنی‌ها می‌افتد.
    \nextgroupplot[
        ymode=log,
        ymin=0.005, ymax=1.0,
        ylabel={\rl{هزینهٔ آموزش هر دوره}},
    ]
    \addplot[draw=thControl, mark=none] table[x=step, y=collapsed, col sep=space]
        {figures/curves/train_avg_loss_per_epoch.csv};
    \addplot[draw=thIdeal, mark=none] table[x=step, y=fixed, col sep=space]
        {figures/curves/train_avg_loss_per_epoch.csv};
    \node[thnote, anchor=west, text=thControl] at (axis cs:1.5,0.155)
        {\rl{اجرای فروریخته --- هدف: عمق مستقیم}};
    \node[thnote, anchor=west, text=thIdeal] at (axis cs:1.5,0.031)
        {\rl{اجرای اصلاح‌شده --- هدف: عمق معکوس}};

    \nextgroupplot[
        ymin=0, ymax=0.46,
        ylabel={\rl{\lr{AbsRel} اعتبارسنجی}},
        xlabel={\rl{دورهٔ آموزش}},
    ]
    \addplot[draw=thControl, mark=none] table[x=step, y=collapsed, col sep=space]
        {figures/curves/eval_abs_rel.csv};
    \addplot[draw=thIdeal, mark=none] table[x=step, y=fixed, col sep=space]
        {figures/curves/eval_abs_rel.csv};
    \node[thnote, anchor=west, text=thControl] at (axis cs:1.5,0.415)
        {\rl{اجرای فروریخته --- منجمد روی $0.3465$ از دورهٔ نخست}};
    \node[thnote, anchor=west, text=thIdeal] at (axis cs:1.5,0.135)
        {\rl{اجرای اصلاح‌شده --- نزول عادی تا $0.055$}};

    \end{groupplot}
    \end{tikzpicture}
    \caption[امضای فروریزش خاموش در منحنی‌های آموزش]{امضای \emph{فروریزش
    خاموش} توضیح‌داده‌شده در بخش~\ref{subsec:disparity_target_space}، در دو
    اجرای کامل ۴۰ دوره‌ای \lr{ViT-L} که تنها در فضای هدف تفاوت دارند.
    \textbf{بالا:} هزینهٔ آموزش اجرای فروریخته (قرمز) هموار و نزولی است ---
    از \fadecimal{۰٫۴۴۲} به \fadecimal{۰٫۴۲۹} --- و اگر تنها همین منحنی پایش
    شود، آموزش کاملاً سالم به نظر می‌رسد؛ زیرا با خروجی ثابت، مقدار هزینه
    صرفاً تابع آماره‌های هر دسته است. \textbf{پایین:} معیار اعتبارسنجی همان
    اجرا از \emph{دورهٔ نخست} روی \fadecimal{۰٫۳۴۶۵} منجمد می‌ماند و تا پایان
    آموزش حتی در رقم چهارم اعشار تغییر نمی‌کند. نکتهٔ گویا آنکه در این حالت،
    مقادیر اعتبارسنجی دو بازوی \proposedmodel{} و \controlmodel{}
    \emph{بیت‌به‌بیت} یکسان است، چون خروجی هر دو شبکه ثابت صفر شده و دیگر به
    ورودی --- و در نتیجه به $K$ --- وابسته نیست. پس از انتقال هدف به فضای عمق
    معکوس (سبز)، هر دو منحنی رفتار عادی می‌یابند. این شکل نشان می‌دهد که چرا
    پایش هزینهٔ آموزش به‌تنهایی برای تشخیص سلامت آموزش کافی نیست.}
    \label{fig:silent_collapse_curves}
\end{figure}

% =============================================================================

\begin{figure}[htbp]
    \centering
    \begin{tikzpicture}
    \begin{groupplot}[
        group style={group size=1 by 2, vertical sep=1.05cm, x descriptions at=edge bottom},
        thesisplot,
        width=0.86\textwidth, height=5.0cm,
        xmin=-1, xmax=40,
    ]

    % منحنی‌ها مستقیماً برچسب‌گذاری می‌شوند و راهنمای جعبه‌ای حذف شده است، چون
    % در این دو نمودار جعبهٔ راهنما ناگزیر روی خودِ منحنی‌ها می‌افتد.
    \nextgroupplot[
        ymode=log,
        ymin=0.005, ymax=1.0,
        ylabel={\rl{هزینهٔ آموزش هر دوره}},
    ]
    \addplot[draw=thControl, mark=none] table[x=step, y=collapsed, col sep=space]
        {figures/curves/train_avg_loss_per_epoch.csv};
    \addplot[draw=thIdeal, mark=none] table[x=step, y=fixed, col sep=space]
        {figures/curves/train_avg_loss_per_epoch.csv};
    \node[thnote, anchor=west, text=thControl] at (axis cs:1.5,0.155)
        {\rl{اجرای فروریخته --- هدف: عمق مستقیم}};
    \node[thnote, anchor=west, text=thIdeal] at (axis cs:1.5,0.031)
        {\rl{اجرای اصلاح‌شده --- هدف: عمق معکوس}};

    \nextgroupplot[
        ymin=0, ymax=0.46,
        ylabel={\rl{\lr{AbsRel} اعتبارسنجی}},
        xlabel={\rl{دورهٔ آموزش}},
    ]
    \addplot[draw=thControl, mark=none] table[x=step, y=collapsed, col sep=space]
        {figures/curves/eval_abs_rel.csv};
    \addplot[draw=thIdeal, mark=none] table[x=step, y=fixed, col sep=space]
        {figures/curves/eval_abs_rel.csv};
    \node[thnote, anchor=west, text=thControl] at (axis cs:1.5,0.415)
        {\rl{اجرای فروریخته --- منجمد روی $0.3465$ از دورهٔ نخست}};
    \node[thnote, anchor=west, text=thIdeal] at (axis cs:1.5,0.135)
        {\rl{اجرای اصلاح‌شده --- نزول عادی تا $0.055$}};

    \end{groupplot}
    \end{tikzpicture}
    \caption[امضای فروریزش خاموش در منحنی‌های آموزش]{امضای \emph{فروریزش
    خاموش} توضیح‌داده‌شده در بخش~\ref{subsec:disparity_target_space}، در دو
    اجرای کامل ۴۰ دوره‌ای \lr{ViT-L} که تنها در فضای هدف تفاوت دارند.
    \textbf{بالا:} هزینهٔ آموزش اجرای فروریخته (قرمز) هموار و نزولی است ---
    از \fadecimal{۰٫۴۴۲} به \fadecimal{۰٫۴۲۹} --- و اگر تنها همین منحنی پایش
    شود، آموزش کاملاً سالم به نظر می‌رسد؛ زیرا با خروجی ثابت، مقدار هزینه
    صرفاً تابع آماره‌های هر دسته است. \textbf{پایین:} معیار اعتبارسنجی همان
    اجرا از \emph{دورهٔ نخست} روی \fadecimal{۰٫۳۴۶۵} منجمد می‌ماند و تا پایان
    آموزش حتی در رقم چهارم اعشار تغییر نمی‌کند. نکتهٔ گویا آنکه در این حالت،
    مقادیر اعتبارسنجی دو بازوی \proposedmodel{} و \controlmodel{}
    \emph{بیت‌به‌بیت} یکسان است، چون خروجی هر دو شبکه ثابت صفر شده و دیگر به
    ورودی --- و در نتیجه به $K$ --- وابسته نیست. پس از انتقال هدف به فضای عمق
    معکوس (سبز)، هر دو منحنی رفتار عادی می‌یابند. این شکل نشان می‌دهد که چرا
    پایش هزینهٔ آموزش به‌تنهایی برای تشخیص سلامت آموزش کافی نیست.}
    \label{fig:silent_collapse_curves}
\end{figure}

% =============================================================================

\begin{figure}[htbp]
    \centering
    \begin{tikzpicture}
    \begin{groupplot}[
        group style={group size=1 by 2, vertical sep=1.05cm, x descriptions at=edge bottom},
        thesisplot,
        width=0.86\textwidth, height=5.0cm,
        xmin=-1, xmax=40,
    ]

    % منحنی‌ها مستقیماً برچسب‌گذاری می‌شوند و راهنمای جعبه‌ای حذف شده است، چون
    % در این دو نمودار جعبهٔ راهنما ناگزیر روی خودِ منحنی‌ها می‌افتد.
    \nextgroupplot[
        ymode=log,
        ymin=0.005, ymax=1.0,
        ylabel={\rl{هزینهٔ آموزش هر دوره}},
    ]
    \addplot[draw=thControl, mark=none] table[x=step, y=collapsed, col sep=space]
        {figures/curves/train_avg_loss_per_epoch.csv};
    \addplot[draw=thIdeal, mark=none] table[x=step, y=fixed, col sep=space]
        {figures/curves/train_avg_loss_per_epoch.csv};
    \node[thnote, anchor=west, text=thControl] at (axis cs:1.5,0.155)
        {\rl{اجرای فروریخته --- هدف: عمق مستقیم}};
    \node[thnote, anchor=west, text=thIdeal] at (axis cs:1.5,0.031)
        {\rl{اجرای اصلاح‌شده --- هدف: عمق معکوس}};

    \nextgroupplot[
        ymin=0, ymax=0.46,
        ylabel={\rl{\lr{AbsRel} اعتبارسنجی}},
        xlabel={\rl{دورهٔ آموزش}},
    ]
    \addplot[draw=thControl, mark=none] table[x=step, y=collapsed, col sep=space]
        {figures/curves/eval_abs_rel.csv};
    \addplot[draw=thIdeal, mark=none] table[x=step, y=fixed, col sep=space]
        {figures/curves/eval_abs_rel.csv};
    \node[thnote, anchor=west, text=thControl] at (axis cs:1.5,0.415)
        {\rl{اجرای فروریخته --- منجمد روی $0.3465$ از دورهٔ نخست}};
    \node[thnote, anchor=west, text=thIdeal] at (axis cs:1.5,0.135)
        {\rl{اجرای اصلاح‌شده --- نزول عادی تا $0.055$}};

    \end{groupplot}
    \end{tikzpicture}
    \caption[امضای فروریزش خاموش در منحنی‌های آموزش]{امضای \emph{فروریزش
    خاموش} توضیح‌داده‌شده در بخش~\ref{subsec:disparity_target_space}، در دو
    اجرای کامل ۴۰ دوره‌ای \lr{ViT-L} که تنها در فضای هدف تفاوت دارند.
    \textbf{بالا:} هزینهٔ آموزش اجرای فروریخته (قرمز) هموار و نزولی است ---
    از \fadecimal{۰٫۴۴۲} به \fadecimal{۰٫۴۲۹} --- و اگر تنها همین منحنی پایش
    شود، آموزش کاملاً سالم به نظر می‌رسد؛ زیرا با خروجی ثابت، مقدار هزینه
    صرفاً تابع آماره‌های هر دسته است. \textbf{پایین:} معیار اعتبارسنجی همان
    اجرا از \emph{دورهٔ نخست} روی \fadecimal{۰٫۳۴۶۵} منجمد می‌ماند و تا پایان
    آموزش حتی در رقم چهارم اعشار تغییر نمی‌کند. نکتهٔ گویا آنکه در این حالت،
    مقادیر اعتبارسنجی دو بازوی \proposedmodel{} و \controlmodel{}
    \emph{بیت‌به‌بیت} یکسان است، چون خروجی هر دو شبکه ثابت صفر شده و دیگر به
    ورودی --- و در نتیجه به $K$ --- وابسته نیست. پس از انتقال هدف به فضای عمق
    معکوس (سبز)، هر دو منحنی رفتار عادی می‌یابند. این شکل نشان می‌دهد که چرا
    پایش هزینهٔ آموزش به‌تنهایی برای تشخیص سلامت آموزش کافی نیست.}
    \label{fig:silent_collapse_curves}
\end{figure}
